% Root and DE independent audits passed on 2026-11-15.
\begin{statement}{C.33}{Proved}{An arithmetic-trivial character of order thirteen}
There is a continuous surjective homomorphism
\[
 \epsilon_{13}:\widehat{\mathrm{GT}}_0\longrightarrow\mathbb F_{13}
\]
which vanishes on the absolute Galois image and is nonzero on the
multiplier-one subgroup.  It does not factor through the actual
metabelian image of B.7.  Consequently the abelianization of the full
arithmetic-normal quotient has an odd-order quotient; it is neither
an elementary abelian two-group nor a pro-two group.
\end{statement}
\begin{proof}
We first justify the actual local image and its integral filtration.
Let $F_{13}$ be the free pro-$13$ group on $x,y$.  Denote by
$\mathcal G$ the image of the full gentle group in its normalized
automorphism group, and by $\mathcal U$ the kernel of the local
multiplier $\lambda_{13}:\mathcal G\to\mathbb Z_{13}^{\times}$.
For $n\geq2$, let $\mathcal U_n$ consist of the elements whose
normalized conjugator belongs to $\gamma_nF_{13}$.

Put $z=-x-y$ in the integral free Lie algebra $L$ on $x,y$.
For a coefficient ring $R$, let $S_n(R)$ be the degree-$n$ solutions of
\[
\begin{split}
 q(x,y)+q(y,x)&=0,\\
 [x,q(x,y)]&=[x+y,q(-x-y,y)],\\
 q(x,y)+q(y,z)+q(z,x)&=0.
\end{split}                                                    \tag{C33.1}
\]
These are exactly the normalized special outer derivations fixed
by the geometric $S_3$ action.  To check the assertion, write
$D_q(x)=0$, $D_q(y)=[y,q]$.  The second equation says that
$D_q(z)=[z,q(x,y)-q(z,y)]$.  The transposition of $x,y$
changes $D_q$ by $\operatorname{ad}(q)$, and the cycle of $x,y,z$
does the same precisely when the second equation holds.
Its third iterate gives the third equation: the sum of the three
inner conjugators must be central, and the free Lie algebra has
zero center.  Conversely these computations give both required
outer invariances.  No degree-$n$ peripheral centralizer ambiguity
occurs for $n\geq2$.

The following equality concerns the actual completed groups:
\[
 \mathcal U_n/\mathcal U_{n+1}\simeq S_n(\mathbb Z_{13}).
                                                               \tag{C33.2}
\]
Necessity follows by taking the first nonzero Magnus terms of the
peripheral and $S_3$ equations.  For surjectivity, apply the general
integral argument inside P.2 to any $q$ in (C33.1), not just to its
particular degree-seven example.  Lift its three peripheral
conjugators to $F_{13}$.  At each subsequent degree $m$, the error
in their product relation can be changed by
$(u,v)\mapsto[u,x]+[v,y]$ in $L_m$; this map is integrally onto
by Jacobi induction on Lie words.  The corrections converge,
give an automorphism by the pro-$13$ basis theorem, and do not
change the original homogeneous derivation.  Coprime complement
conjugacy for the $S_3$ action then removes all discrepancies in
strictly higher degree.  Thus this is a complete pro-$13$
gentle automorphism.  The relative pronilpotent construction of
P.4 lifts that entire solution to the full gentle group, using
the identity at all other prime components and retaining the
cusp conjugator.  It has global multiplier one.  This proves
(C33.2) for $\mathcal U$ as defined, rather than for a group
of merely truncated solutions.

The exact integral computation gives the following ranks:
\[
\begin{array}{c|rrrrrrrrrrr}
 n&2&3&4&5&6&7&8&9&10&11&12\\ \hline
 \operatorname{rank}S_n&0&1&0&1&0&2&1&4&2&9&7\\
 \operatorname{rank}[S,S]_n&0&0&0&0&0&0&1&0&2&1&6 .
\end{array}                                                    \tag{C33.3}
\]
The two programs \texttt{pbc\_gentle\_even.py} and
\texttt{pbc\_gentle\_even\_fast.py} independently evaluate the
substitutions by associative expansion and recursive Lyndon
brackets.  Their complete mathematical outputs agree.
For every weight, every displayed integral kernel vector is
checked against every original equation row.  An independent
rank lower bound modulo $65521$ equals the complementary
dimension, proving the rational ranks.  The further certificate
\texttt{pbc\_gentle\_weight12\_integral.py} checks a unit minor
modulo $13$ for each kernel basis.  Since the full integral
kernel is saturated, these are bases over $\mathbb Z_{13}$.
It also computes the Ihara bracket
\[
 \{f,g\}=[f,g]+D_f(g)-D_g(f)
\]
and checks every bracket against all the defining rows.
In the retained seven-element basis of $S_{12}$, the primitive
integral functional
\[
             \ell=(-4,1,8,-4,-17,14,84)                         \tag{C33.4}
\]
annihilates every degree-twelve bracket.  Multiplying its reduction
by $3$ gives $(1,3,11,1,1,3,5)$, so the first basis vector has
value $1$.  Thus this is an integral witness, with no inference
from rational bracket membership modulo $13$.

Now put $U=\mathcal U/\mathcal U_{13}$.
The zero degree-two space in (C33.3) implies
$\mathcal U=\mathcal U_3$.  The commutator filtration gives
$[U_a,U_b]\subseteq U_{a+b}$, so $U$ has nilpotency class at most
four.  It is a finitely generated torsion-free pro-$13$ group:
its finite filtration has free finite-rank $\mathbb Z_{13}$
quotients.  Its Lazard Lie lattice
$\mathfrak u=\log U$ is defined integrally, preserves this
filtration and its associated graded brackets, and is functorial
for all automorphisms.
For a direct denominator check, use the action on the integral
Magnus algebra modulo augmentation degree $14$.
For $a\in U$, $a-1$ raises augmentation degree by at least three.
The operator logarithm and exponential therefore stop after
four terms; their denominators are units in $\mathbb Z_{13}$.
The Lazard correspondence applies to the actual closed subgroup
of these operators, since its class is less than $13$.
No logarithm of a free-group generator, and no class-thirteen
Lazard assertion about the free group, is used here.
See \cite[Theorem 1]{PBC_Garaialde2014} for the precise
finitely generated pro-prime correspondence, including its
subgroup and automorphism assertions.

Set $G=\mathcal G/\mathcal U_{13}$, so that
\[
       1\longrightarrow U\longrightarrow G
        \xrightarrow{\lambda_{13}}\mathbb Z_{13}^{\times}
        \longrightarrow1 .                                    \tag{C33.5}
\]
Its multiplier is onto because the Galois multiplier is onto.
In the local Galois image, the inverse image of the Teichmuller
group $\mu_{12}$ is an extension of $\mu_{12}$ by a pro-$13$
group.  The coprime complement theorem and compactness give an
element $t$ of order $12$ with primitive multiplier
$\omega\in\mu_{12}$.
Conjugation by $t$ acts on
$\operatorname{gr}_n\mathfrak u=S_n(\mathbb Z_{13})$
as $\omega^n$.  The ten possible weights $3,\ldots,12$
have distinct residues modulo $12$.  Hence the integral
idempotents
\[
 e_n=\frac1{12}\sum_{j=0}^{11}\omega^{-nj}\operatorname{Ad}(t)^j
\]
split the filtered lattice, and
$e_n\mathfrak u\to\operatorname{gr}_n\mathfrak u$ is an
isomorphism.  There can be no kernel or cokernel in this map:
their successive graded pieces would have to carry both
the eigencharacter $\omega^n$ and a distinct $\omega^m$.
The Lie bracket respects this splitting.  If $a+b\leq12$,
the bracket of the $a$ and $b$ summands is exactly the
degree-$a+b$ Ihara bracket; if $a+b>12$ it is zero.
This removes the possibility of hidden higher filtered
commutator relations.

Extend $\ell$ by zero on all eigensummands except weight $12$.
It annihilates $[\mathfrak u,\mathfrak u]$.  Its reduction,
normalized as after (C33.4), therefore gives a nonzero
continuous homomorphism
\[
                         \theta:U\longrightarrow\mathbb F_{13}.
                                                               \tag{C33.6}
\]
This character is invariant under the entire $G$.
Indeed conjugation on
$\mathfrak u/[\mathfrak u,\mathfrak u]$ factors through
$G/U=\mathbb Z_{13}^{\times}$ and hence commutes with $t$.
It preserves the weight-twelve eigensummand.  That summand
is the top filtered piece, where an element of multiplier
$a$ acts exactly by $a^{12}$, which is $1$ modulo $13$.
This proves invariance without assuming that a chosen wild
multiplier lift commutes with $t$ inside $G$.

There is no extension obstruction:
\[
 H^2_{\mathrm{cts}}(\mathbb Z_{13}^{\times},\mathbb F_{13})=0.
                                                               \tag{C33.7}
\]
The prime-to-$13$ factor $\mu_{12}$ has vanishing positive
cohomology by averaging; the procyclic factor
$1+13\mathbb Z_{13}\simeq\mathbb Z_{13}$ has
$13$-cohomological dimension one.  Equivalently, the
inflation maps in degree two along its finite cyclic
quotients multiply by $13$, so their direct limit is zero.
The continuous five-term sequence of (C33.5) now extends
$\theta$ to a continuous homomorphism
$\widetilde\theta:G\to\mathbb F_{13}$.

Finally remove its arithmetic restriction integrally.
The normalized pro-$13$ action on the tripod, with its
standard tangential basepoint, is unramified away $13$;
this is stated for the actual pro-prime automorphism
representation in \cite[\S8.5, immediately before
Theorem 8.7]{A_HainMatsumoto2003}, taking $N=1$.
Thus $\widetilde\theta|_{G_{\mathbb Q}}$ factors through
$G_{\mathbb Q,\{13\}}$.
By Kronecker--Weber, the maximal abelian extension of
$\mathbb Q$ unramified away $13$ is
$\mathbb Q(\mu_{13^\infty})$; see
\cite[Chapter I, Theorem 4.16 and its proof]{PBC_MilneCFT}.
Consequently
$\operatorname{Hom}_{\mathrm{cts}}
(G_{\mathbb Q,\{13\}},\mathbb F_{13})$ is one-dimensional,
generated by the nonzero cyclotomic character
\[
       c_{13}(\lambda)=\frac{\lambda_{13}^{12}-1}{13}\pmod {13}.
\]
This expression is an ordinary additive character on the
whole multiplier group: multiplication of its two numerators
has cross term divisible by $13^2$.
There is therefore $a\in\mathbb F_{13}$ such that
\[
 \epsilon_{13}=\widetilde\theta\circ\rho_{13}
                  -a\,c_{13}\circ\lambda
                                                               \tag{C33.8}
\]
vanishes on $G_{\mathbb Q}$.  The adjustment is zero on
global multiplier one.  The actual complete lift supplied
by (C33.2) for the first weight-twelve basis vector still
has value $1$, so (C33.8) is surjective.

If it factored through the actual metabelian image, it
would factor through that image's arithmetic-normal
quotient.  C.32 proves that quotient has exponent two,
which has no nonzero homomorphism to $\mathbb F_{13}$.
This contradiction proves the last assertion.
\end{proof}
